\section{Related work and scope of the contribution}
\label{sec:related}

\paragraph{Equilibration and numerical reliability.}
Diagonal scaling is an established numerical intervention. Classical methods include
the logarithmic scaling of Curtis and Reid \cite{curtis1972scaling}; the computational
comparison by Elble and Sahinidis \cite{elble2012scaling} shows that additional
complexity does not necessarily improve simplex performance over simple
equilibration. Klotz \cite{klotz2014numerics} explains how finite precision,
ill-conditioning and formulation choices produce inaccurate or inconsistent LP and
MIP results. These precedents motivate separate measurements of coefficient
magnitudes, numerical reliability and computational effort. Neither the validity of
positive row scaling nor the observation that conditioning and runtime can disagree
is a contribution of the present work.

Optimal diagonal preconditioning is a broader problem than coefficient centering.
Qu et al. \cite{qu2024diagonal} formulate condition-number minimization through
quasiconvex and semidefinite optimization, and Demmel \cite{demmel2023block}
provides guarantees for block-Jacobi scaling. The scalar triangular calculation
below is a tractable special case used to explain the coupling-kernel mechanism;
it is not a general optimal-preconditioning algorithm.

Robustness also has a direct antecedent in MIP scaling. Berthold and Hendel
\cite{berthold2021scaling} select between standard and Curtis--Reid scaling in
Xpress by learning an indicator of numerical difficulties, rather than optimizing
runtime alone. Their implementation uses power-of-two factors and their evaluation
includes inconsistencies across seeds. We instead study an external rule with explicit
row-level admissibility conditions. Its conditions concern a declared residual budget
and coefficient window; they do not predict the solver's numerical behavior or
replace its internal scaling.

\paragraph{Scaling at the modeling layer.}
User-defined scaling is supported by modeling environments: Pyomo's
\texttt{ScaleModel} applies factors attached to model components
\cite{pyomoScaleModel}, and GAMS provides equation and variable scale attributes
\cite{gamsScaling}. IDAES goes further by storing factors within the owning blocks,
deriving constraint scales from variable and expression scales, and providing
propagation between submodels \cite{idaesScaling,idaesCustomScaler}. Graph-based
modeling in Plasmo.jl likewise preserves local problems and linking constraints
\cite{jalving2022plasmo}. Thus using generator information before solver invocation
is established practice. In the setting considered here, the additional input is a
declared budget for each original residual. A row's block label or coefficient norm
alone does not determine that budget.

\paragraph{Tolerances and application units.}
Gurobi's numerical guide explains that its feasibility and integrality tolerances
are absolute, and that very small constraint coefficients may be discarded
\cite{gurobiScaling}. Its modeling guidance explicitly recommends external scaling
that makes tolerance-sized violations negligible in the application
\cite{gurobiExternalScaling}. Tolerance-based constraint scaling also appears
outside MIP: Pritchett, Howell and Grebow \cite{pritchett2017scaling} choose scales
from the ratio between a desired physical accuracy and SNOPT's feasibility
tolerance. The elementary transfer of a residual budget through a positive row
factor therefore has clear precedent.

The construction in Section~\ref{sec:contract} combines this transfer with a
prescribed coefficient window. It characterizes the admissible factors, gives the
constrained geometric-mean choice, and distinguishes two reasons why a row-only
intervention cannot satisfy the requested specifications. Its integer-rounding
extension isolates a further error term that is unaffected by row scaling. These
results organize familiar numerical ingredients into a checkable contract; the
empirical question is whether that contract and its abstention cases improve the
interpretation and reliability of preprocessing decisions. No general superiority
over existing scaling methods follows from the characterization.

\paragraph{Checking results and checking attribution.}
Independent verification of MIP results is a broader task. Cheung, Gleixner and
Steffy \cite{cheung2017certificates} introduce certificates for checking claimed
integer-programming results, while Hoen and Gleixner \cite{hoen2024correctness}
analyze the numerical correctness of individual branch-and-bound decisions.
The exact original-data residual checks used here address primal acceptance;
they do not certify a solver's global lower bound. Optimality is checked separately
by enumeration only on our small stress instances.

The archived experiments provide a complementary attribution study. A change in
row coverage or in the norm used on linking rows can explain a difference initially
ascribed to block information. We therefore compare actual transformations under
matched coverage and kernels, including algebraic simplification and export
identity checks. The resulting conclusions concern the mechanisms and instance
families tested, rather than an impossibility of useful structure-aware scaling.
